\documentclass[letterpaper,landscape]{article}
\usepackage[margin=0.32in]{geometry}
\usepackage[T1]{fontenc}
\usepackage{lmodern,amsmath,amssymb,multicol,xcolor,enumitem,microtype}
\usepackage[hidelinks]{hyperref}
\definecolor{navy}{RGB}{24,49,72}
\definecolor{rulegray}{RGB}{168,177,184}
\pagestyle{empty}
\setlength{\parindent}{0pt}
\setlength{\parskip}{2.2pt}
\setlength{\columnsep}{0.19in}
\setlength{\columnseprule}{0.3pt}
\def\columnseprulecolor{\color{rulegray}}
\setlength{\abovedisplayskip}{3pt}
\setlength{\belowdisplayskip}{3pt}
\setlength{\abovedisplayshortskip}{2pt}
\setlength{\belowdisplayshortskip}{2pt}
\DeclareMathOperator{\E}{E}
\DeclareMathOperator{\Var}{Var}
\DeclareMathOperator{\Cov}{Cov}
\newcommand{\dd}{\,d}
\newcommand{\head}[1]{\vspace{3pt}{\color{navy}\bfseries\fontsize{10}{11}\selectfont #1}\par\vspace{1pt}\hrule\vspace{3pt}}
\newcommand{\mini}[1]{{\bfseries #1}}
\newcommand{\cmd}[1]{{\ttfamily\fontsize{7.6}{8.7}\selectfont #1}}
\newcommand{\topline}[2]{{\color{navy}\bfseries\fontsize{15}{16}\selectfont STA 431\quad Midterm II}\hfill{\bfseries #1}\par\vspace{2pt}{\fontsize{8}{9}\selectfont #2}\par\vspace{3pt}\hrule\vspace{3pt}}
\begin{document}
\fontsize{8.3}{9.35}\selectfont
\topline{FRONT: DISTRIBUTIONS + MOMENTS}{Handwritten-copy master\quad\textbullet\quad US Letter, landscape\quad\textbullet\quad Three columns on each side}
\begin{multicols*}{3}
\head{1\quad Joint, marginal, conditional}
\mini{Normalize first.} PMF: $p(x,y)\ge0$, $\sum_x\sum_y p(x,y)=1$. PDF: $f(x,y)\ge0$, $\iint_S f=1$. If $f=cq$, then $c=1/\iint_Sq$ (use sums for PMFs). Write the support $S$; zero elsewhere.

\mini{Marginal = eliminate the other variable.}
\[
 p_X(x)=\sum_y p(x,y),\qquad
 f_X(x)=\int_{y:(x,y)\in S}f(x,y)\dd y.
\]
Interchange $x,y$ for the other marginal. A marginal's support is the projection of $S$ onto its axis.

\mini{Conditional = joint / conditioning marginal.}
\[
 p_{X\mid Y}(x\mid y)=\frac{p(x,y)}{p_Y(y)},\qquad
 f_{X\mid Y}(x\mid y)=\frac{f(x,y)}{f_Y(y)}.
\]
Use only where the denominator is positive; retain the slice of $S$ at that $y$. Sum/integrate to 1 over $x$.
\[
 f(x,y)=f_{X\mid Y}(x\mid y)f_Y(y).
\]
A conditional alone or two marginals generally do \emph{not} determine the joint. Conditional plus its conditioning marginal does (almost everywhere).

\mini{Joint CDF.} $F(x,y)=P(X\le x,Y\le y)$.
\[
 F(x,y)=\sum_{u\le x}\sum_{v\le y}p(u,v)
\]
\[
 F(x,y)=\int_{-\infty}^{x}\int_{-\infty}^{y}f(u,v)\dd v\dd u.
\]
Intersect this rectangle with $S$ before integrating.
\[
 F_X(x)=\lim_{y\to\infty}F(x,y),\quad
 F_Y(y)=\lim_{x\to\infty}F(x,y).
\]
For an absolutely continuous joint law,
$f(x,y)=\partial^2F/\partial x\partial y$ almost everywhere.
For integer-valued variables,
\[
\begin{split}
p(x,y)={}&F(x,y)-F(x-1,y)\\
&-F(x,y-1)+F(x-1,y-1).
\end{split}
\]
\mini{Rectangle from any CDF:}
\[
\begin{split}
P(a<X\le b,c<Y\le d)
={}&F(b,d)-F(a,d)\\
&-F(b,c)+F(a,c).
\end{split}
\]
For discrete inclusive lower endpoints use left limits (or $a-1,c-1$ on an integer lattice). Continuous endpoints have zero mass.

\columnbreak
\head{2\quad Events, bounds, independence}
\mini{Region method.} Translate words $\to$ event $A$; draw $S\cap A$; select order; find outer range; then inner slice limits.
\[
 P((X,Y)\in A)=\sum_{(x,y)\in S\cap A}p(x,y)
 =\iint_{S\cap A} f(x,y)\dd x\dd y.
\]
An empty slice contributes zero; split where a boundary changes. The inner limits may depend on the outer variable; the outer limits cannot depend on the inner variable.

\mini{Common regions.} On $0<x<1,0<y<2$:
\[
 P(X<Y)=\int_0^1\int_x^2 f(x,y)\dd y\dd x.
\]
Equivalently split at $y=1$:
\[
 \int_0^1\int_0^y f\dd x\dd y+
 \int_1^2\int_0^1 f\dd x\dd y.
\]
On $x>0,y>0,x+y<1$:
\[
 \iint_S f=\int_0^1\int_0^{1-x}f\dd y\dd x
 =\int_0^1\int_0^{1-y}f\dd x\dd y.
\]
For $P(X+Y\le a)$, intersect $y\le a-x$ with the support. For $P(X>Y)$, use $y<x$. In a PMF list actual allowed pairs; do not integrate.

\mini{Words $\to$ notation.}
\begin{itemize}[leftmargin=10pt,itemsep=1pt,topsep=1pt]
\item Both lifetimes exceed $t$: $X>t,Y>t$.
\item Series device fails if either fails:
$T=\min(X,Y)$; $P(T\le t)=1-P(X>t,Y>t)$.
\item Parallel device fails only if both fail:
$T=\max(X,Y)$; $P(T\le t)=P(X\le t,Y\le t)$.
\item At least $a$ given $Y=y$: integrate $f_{X\mid Y}$ over $x\ge a$ in its slice.
\end{itemize}
For an event $B$ of positive probability,
$P(A\mid B)=P(A\cap B)/P(B)$.
For continuous conditioning $Y=y$, use the conditional density; $P(Y=y)=0$.

\mini{Independence.} For all $x,y$,
\[
 F(x,y)=F_X(x)F_Y(y).
\]
Equivalent: $p(x,y)=p_X(x)p_Y(y)$ at every pair; or $f(x,y)=f_X(x)f_Y(y)$ almost everywhere, \emph{including zero regions}. One failing pair (PMF) or a positive-area failure (PDF) disproves it. Factorable formula on triangular support is insufficient. Independence makes the conditional equal its marginal and permits factorization of event probabilities and expectations of products.

\columnbreak
\head{3\quad Expectation, moments, MGFs}
\mini{Expectation / LOTUS.}
\[
\E[g(X)]=\sum_x g(x)p_X(x)
\quad\text{or}\quad\int g(x)f_X(x)\dd x.
\]
For $g(X,Y)$ use $\sum\sum g p$ or $\iint_S g f$.
Finite expectation requires $\E|g(X)|<\infty$.
\[
\E[aX+bY+c]=a\E X+b\E Y+c.
\]
Linearity needs no independence. In general $\E[g(X)]\ne g(\E X)$. For an indicator, $\E[1_A]=P(A)$. If $g\ge0$, then $\E g\ge0$.

\mini{Raw and central moments.}
\[
 \mu_r'=\E[X^r],\quad\mu=\mu_1',\quad
 \mu_r=\E[(X-\mu)^r].
\]
\[
 \mu_0'=1,\quad\mu_1=0,\quad
 \Var X=\mu_2=\mu_2'-\mu^2.
\]
\[
 \Var(aX+b)=a^2\Var X,\quad
 \sigma_{aX+b}=|a|\sigma_X.
\]
For $Y=g(X)$, compute $\E g$ and $\E g^2$; then
$\Var Y=\E[g(X)^2]-(\E g(X))^2$.
For $Y=X^2$, variance requires $\E X^4$.
If density/support are symmetric about zero, existing odd moments vanish; even moments need integration.

\mini{MGF.} $M_X(t)=\E[e^{tX}]$.
\[
 M_X(t)=\sum_x e^{tx}p_X(x)
 \quad\text{or}\quad\int e^{tx}f_X(x)\dd x.
\]
State the set of $t$ where it is finite. The usual moment and uniqueness results require finiteness on an open interval containing $0$.
\[
 M(0)=1,\quad M^{(r)}(0)=\E X^r,
 \quad\Var X=M''(0)-[M'(0)]^2.
\]
Equal MGFs on a neighborhood of zero imply equal distributions. All moments existing does not guarantee such an MGF.
\[
 M_{aX+b}(t)=e^{bt}M_X(at).
\]
For independent $X_i$,
\[
 M_{\sum a_iX_i+b}(t)=e^{bt}\prod_i M_{X_i}(a_it).
\]
Intersect the transformed domains. Without independence, do not multiply marginal MGFs.

\mini{Fast templates.}
\[
\begin{array}{c|c|c}
\text{Law/form}&M(t)&\E X,\ \Var X\\\hline
\lambda e^{-\lambda x},\ x>0&\frac\lambda{\lambda-t},\ t<\lambda&1/\lambda,\ 1/\lambda^2\\
e^{\alpha t+\beta t^2}&\text{all }t&\alpha,\ 2\beta
\end{array}
\]
$\int_0^\infty x^n e^{-ax}\dd x=n!/a^{n+1}$ for $a>0$, integer $n\ge0$.
\end{multicols*}
\newpage
\topline{BACK: COVARIANCE + CONDITIONING + TI-84}{Use the formulas to set up exact work; use numerical commands to check arithmetic and one-dimensional integrals.}
\begin{multicols*}{3}
\head{4\quad Product moments + linear combinations}
\mini{Product moments.}
\[
 \mu'_{rs}=\E[X^rY^s],\qquad
 \mu_{rs}=\E[(X-\mu_X)^r(Y-\mu_Y)^s].
\]
Compute from the joint law, not by multiplying marginals unless independence is justified.
\[
 \Cov(X,Y)=\E[XY]-\mu_X\mu_Y,
 \qquad \rho=\frac{\Cov(X,Y)}{\sigma_X\sigma_Y}.
\]
Correlation needs positive finite variances; $-1\le\rho\le1$.
$|\Cov(X,Y)|\le\sigma_X\sigma_Y$.
\[
\begin{gathered}
\Cov(X,Y)=\Cov(Y,X),\quad\Cov(X,X)=\Var X,\\
\Cov(aX+b,cY+d)=ac\Cov(X,Y).
\end{gathered}
\]
Independent $\Rightarrow$ zero covariance (finite second moments). Zero covariance $\not\Rightarrow$ independence.
Example: $X$ symmetric on $(-1,1)$, $Y=X^2$;
$\E X=\E X^3=0$, hence $\Cov(X,Y)=0$, yet $Y$ is determined by $X$. To prove dependence use events $A,B$ for which $P(A\cap B)\ne P(A)P(B)$.

\mini{One combination.} $W=c+\sum_i a_iX_i$:
\[
\begin{split}
 \E W&=c+\sum_i a_i\mu_i,\\
 \Var W&=\sum_i a_i^2\sigma_i^2
       +2\sum_{i<j}a_ia_j\sigma_{ij}.
\end{split}
\]
Here $\sigma_{ij}=\Cov(X_i,X_j)$. Write each unordered covariance pair once, times $2$.
\[
\begin{split}
\Var(aX+bY+cZ)={}&a^2\Var X+b^2\Var Y\\
&+c^2\Var Z+2ab\Cov(X,Y)\\
&+2ac\Cov(X,Z)+2bc\Cov(Y,Z).
\end{split}
\]
Negative coefficients change cross-term signs; constants contribute no variance. Uncorrelated variables suffice to remove the cross terms. Independence is stronger.

\mini{Two combinations.} $U=aX+bY+c$,
$V=dX+eY+f$:
\[
\begin{split}
 \Cov(U,V)={}&ad\Var X+be\Var Y\\
 &+(ae+bd)\Cov(X,Y).
\end{split}
\]
More generally $\Cov(\sum_i a_iX_i,\sum_j b_jX_j)
=\sum_i\sum_j a_ib_j\sigma_{ij}$.
Find $\rho_{UV}$ by dividing this covariance by
$\sqrt{\Var U\Var V}$.

\columnbreak
\head{5\quad Conditional moments + workflows}
\mini{Conditional expectation.} Fix the conditioned value, then average over the variable that remains random:
\[
\begin{split}
 m(y)=\E[X\mid Y=y]
 &=\sum_x x\,p_{X\mid Y}(x\mid y)\\
 &\text{or }\int x f_{X\mid Y}(x\mid y)\dd x.
\end{split}
\]
For $\E[g(X)\mid Y=y]$, replace $x$ by $g(x)$.
\[
\begin{split}
 q(y)&=\E[X^2\mid Y=y],\\
 \Var(X\mid Y=y)&=q(y)-[m(y)]^2.
\end{split}
\]
Do not square before taking the conditional second moment. $\E[X\mid Y]$ is a random variable $m(Y)$; $m(y_0)$ is a number.

\mini{Marginal-first recipe.}
\begin{enumerate}[leftmargin=12pt,itemsep=1pt,topsep=1pt]
\item Find $f_Y(y)$ or $p_Y(y)$ with its support.
\item Divide joint by it, retaining the conditional slice.
\item Integrate/sum $x$ and $x^2$ against that law.
\item Evaluate at the specified $y$; subtract $m(y)^2$ for variance.
\end{enumerate}
\mini{Direct ratio shortcut.}
\[
 \E[g(X)\mid Y=y]
 =\frac{\int g(x)f(x,y)\dd x}{f_Y(y)}
\]
(use sums for discrete). Denominator must be positive. Reconstruct marginal moments by averaging over $Y$:
\[
\begin{gathered}
 \E X=\E[\E(X\mid Y)],\\
 \Var X=\E[\Var(X\mid Y)]
          +\Var[\E(X\mid Y)].
\end{gathered}
\]
Independence gives $\E[g(X)\mid Y]=\E g(X)$ when integrable. A constant conditional mean alone does not prove independence.

\mini{Problem-solving order.}
\begin{enumerate}[leftmargin=12pt,itemsep=1pt,topsep=1pt]
\item Identify PMF / PDF / CDF / MGF / supplied moments.
\item Record support, event, and conditioning direction.
\item Normalize if a constant is unknown.
\item Select formula; show the sum/integral and bounds.
\item Keep fractions through the algebra; check the result.
\end{enumerate}
\mini{Checks.} Probabilities in $[0,1]$; PMFs sum to $1$; PDFs integrate to $1$ (a PDF value may exceed $1$); variances nonnegative; $\E X$ inside the range of $X$ when finite; $M(0)=1$; correlation in $[-1,1]$.

\mini{Useful algebra.}
$\sum_{k=1}^{n}k=n(n+1)/2$;
$\sum_{k=1}^{n}k^2=n(n+1)(2n+1)/6$;
$\sum_{k=1}^{n}k^3=[n(n+1)/2]^2$.

\columnbreak
\head{6\quad TI-84 Plus CE: useful checks}
\mini{Find commands.} \texttt{[2nd][0]} opens CATALOG. Select the actual command token. \texttt{MATH}: \texttt{nDeriv(}, \texttt{fnInt(}; \texttt{[2nd][STAT]} LIST: \texttt{OPS}\,$\to$\,\texttt{seq(}, \texttt{MATH}\,$\to$\,\texttt{sum(}. Use \texttt{MODE: CLASSIC} for the linear syntax below; MathPrint instead shows templates.

\mini{Finite PMF: values in $L_1$, probabilities in $L_2$.}
\par\cmd{sum(L2)}\quad (must be $1$)
\par\cmd{sum(L1*L2)}\quad ($\E X$)
\par\cmd{sum(L1\string^2*L2)-sum(L1*L2)\string^2}\quad ($\Var X$)
\par For unnormalized nonnegative weights $w$ in $L_2$, divide each moment sum by \cmd{sum(L2)} first.

\mini{Generate finite sums without lists.}
\par\cmd{sum(seq(expression,K,first,last))}
\par HW5-style $p(x)=(x^2+x)/40$, $x=1,\ldots,4$:
\par\cmd{sum(seq(K*(K\string^2+K)/40,K,1,4))}
\par gives $\E X=13/4$. For $\E X^2$, use $K^2$ instead of the leading $K$. For $\Var(X^2)$, obtain $\E X^4-(\E X^2)^2$.

\mini{Joint PMF as rows of allowed pairs.}
Put all $x$ values in $L_1$, corresponding $y$ in $L_2$, and joint probabilities in $L_3$.
\par\cmd{sum(L1*L2*L3)}\quad ($\E XY$)
\par\cmd{sum((L1>L2)*L3)}\quad ($P(X>Y)$)
\par\cmd{sum((L1+L2<=3)*L3)}\quad ($P(X+Y\le3)$)
\par Use TEST's single $\le$ token for \texttt{<=}; tokens: \texttt{[2nd][MATH]} TEST. Every joint pair needs its own row; do not pair just the distinct marginal values.

\mini{One-dimensional numerical integrals.}
\par\cmd{fnInt(expression,X,lower,upper)}
\par HW5 check: \cmd{fnInt(X\string^2*3/4*(1-X\string^2),X,-1,1)}
\par gives $1/5$. Replacing $X^2$ by $X^4$ gives $3/35$.
For $f(x,y)=x+y$ on the unit square, at $y=1/2$, conditional density is $x+1/2$:
\par\cmd{fnInt(X*(X+1/2),X,0,1)}\quad ($7/12$)
\par\cmd{fnInt(X\string^2*(X+1/2),X,0,1)}\quad ($5/12$)

\mini{MGF first-derivative check.}
\par\cmd{nDeriv(3/(3-X),X,0)}\quad ($\approx1/3$)
\par $X$ here is the MGF's argument $t$. Derive $M''(0)$ by hand; avoid nested numerical derivatives.

\mini{Limits.} No native symbolic algebra. Do the inner integral by hand before \texttt{fnInt}; numerical output does not prove independence or MGF existence. \texttt{1-Var Stats} with weights uses $\sigma_x^2$ for population variance; \texttt{Sx\string^2} and \texttt{variance(} are sample quantities. Prefer weighted sums above. \texttt{MATH}\,$\to$\,$\blacktriangleright$\texttt{Frac} may suggest a fraction; it does not prove exactness.
\end{multicols*}
\end{document}
